
Standard alignment methods optimize for human preferences but may inadvertently neglect user agency---patterns in responses that help users maintain understanding and control. This work proposes BiDPO (Bidirectional DPO), which augments Direct Preference Optimization with an auxiliary agency loss that penalizes responses lacking collaborative markers such as reasoning traces, uncertainty acknowledgment, and engagement cues. Experiments on Mistral-7B with HH-RLHF data reveal a partial validation: the collaboration score is orthogonal to preference labels (r = $-0.026)$, and BiDPO training is numerically stable (loss decreases from 0.883 to 0.892 over 250 steps with zero NaN/Inf occurrences). However, at proof-of-concept scale, BiDPO does not significantly improve generation-time collaboration scores over baseline DPO (+0.54\%, p = 0.247, Cohen's d = 0.016). This negative result highlights a gap between training-time auxiliary signals and generation-time behavioral change---an important consideration for multi-objective alignment research. Full-scale training, $\lambda$ sweeps, and learned collaboration scores are identified as directions for future work.

